\chapter{Conclusions and future work}
\label{chap:conclusions_and_future_work}

\section{Conclusions}
\label{sec:conclusions}

The supplied proteomics matrix contains a strong, internally reproducible difference between tumour and matched non-tumour tissue. This conclusion is supported by paired abundance, paired detection, pathway organisation, patient-grouped predictive modelling, paired-label permutation, multiverse analysis and exact patient deletion.

The analysis identifies 614 stable abundance candidates, 1,827 stable detection patterns, 271 pathways supported by two internal views, 165 high-priority integrated candidates, and a locked prospective hand-off of 20 targets. It finds no stable patient subtype and no defensible compact predictive panel. The latter negative findings are important because they prevent attractive but unsupported conclusions.

The work demonstrates substantial data-science and AI content while retaining biological and clinical humility. Its main product is not a single classifier or protein list, but a reproducible evidence system that exposes how each claim depends on source rows, analytical choices, patient influence and unresolved limitations.

\section{Immediate computational priority}
\label{sec:immediate_computational_priority}

The immediate priority is to implement Phase 2 exactly as prespecified. It should benchmark missingness and PCA strategies with patient-level artificial masking, paired-effect preservation and stability criteria. After Phase 2, Phases 3–7 should be regenerated without changing frozen thresholds in response to the new results. Differences must be reported as sensitivity evidence.

\section{Required future experimental work}
\label{sec:required_future_experimental_work}

True validation requires new information and therefore lies outside the current data-only scope. A future study should:

\begin{enumerate}
\item recruit patients absent from the discovery workbook;
\item define oral site, pathology, intended-use population and reference standard;
\item include appropriate healthy, benign, inflammatory, premalignant or other-cancer groups if diagnostic specificity is claimed;
\item use orthogonal or explicitly validated quantitative assays;
\item record collection, preservation, batch, run order, technical controls and detection limits;
\item blind laboratory processing and adjudication to candidate rank;
\item lock endpoints, candidate membership and multiplicity before outcome inspection;
\item report failed measurements and exclusions; and
\item assess scientific validity, analytical performance and clinical performance separately.
\end{enumerate}

The Phase 7 power grids should be updated with assay-specific variance, discordance and attrition assumptions rather than the selected discovery effects.
